\documentclass[10pt]{article}
\usepackage[a4paper,top=1.0cm,bottom=0.9cm,left=1.4cm,right=1.4cm]{geometry}
\usepackage[utf8]{inputenc}
\usepackage[T1]{fontenc}
\usepackage{booktabs,enumitem,microtype,hyperref}
\usepackage{titlesec,multicol}
\titlespacing*{\section}{0pt}{4pt}{1pt}
\titleformat*{\section}{\normalsize\bfseries}
\setlist{itemsep=0pt,topsep=1pt,leftmargin=12pt}
\setlength{\parskip}{2pt}\setlength{\parindent}{0pt}
\newcommand{\file}[1]{{\small\texttt{#1}}}
\pagestyle{empty}
\begin{document}
{\large\bfseries Research plan: does the ranking of game-solving RL methods survive stronger tuning, from small exact games to a large market?}\\[1pt]
{\small Figgie half of the joint study with the tsinghua-rl IIG project (Rich), reviewed by him. Written 16 Nov 2026. Course 80245013, final report due 8 Jan 2027. Nothing here has been launched.}

\section{The question, and why this one}
\emph{When random search, CARBS or PROTEIN tune PPO and MMD, with AdamW or Muon, do the rankings of algorithms and memory architectures measured under a default recipe survive, and does the tuning budget a fair ranking needs on small exact-exploitability games carry over to a 4-player hidden-information market (Figgie)?}

We recommend the merge of options (a) and (b), with Figgie as the large-game half. The Figgie memory result alone is too thin to publish: the minGRU gain (58.0\% vs 54.0\% probe, 2 seeds) shrank to 56.5\% vs 53.4\% with seed 2 and is confounded by policy, since each agent is probed on its own games. Rankings at one hand-picked recipe are what Rudolph et al.~\cite{rudolph} warn about, and Rich's re-tuning confirms it on our stack (PTTT, 10 fresh seeds: published MMD 0.171 and PPO 0.168 against re-tuned 0.130 and 0.126, $p=0.0005$, and the tuned methods are not separable, Holm $p=0.21$/$0.43$). Rich runs the factorial (random search, CARBS, PROTEIN $\times$ AdamW, Muon) on Kuhn, Leduc and Rudolph's games with exact exploitability and a CFR+ floor. We run the same factorial on Figgie, where exact exploitability is out of reach, and add the axis Figgie alone can test: memory architecture, scored against an exact Bayesian belief.

\section{Rival hypotheses, and what would be new}
H0, tuning dominates: after strong tuning the arms (PPO vs MMD, none vs minGRU vs GRU) tie within noise, Rudolph's result extended to a large game.
H1, rankings are stable: tuning lifts every arm about equally, so the default-recipe ranking predicts the tuned one.
H2, rankings flip: arms differ in tuning sensitivity~\cite{adkins}, so the winner depends on budget. Memory cores, with more interacting hyperparameters, gain most, and the midterm's ``memory does not clearly help'' is a tuning artefact.
H3, the optimizer matters: Muon's orthogonalised updates help recurrent matrices even at 128 width (against: no effect at this scale).
Prior work and the gap. Rudolph et al.~\cite{rudolph} used random search on small 2-player games only. Eimer et al.~\cite{eimer} compared tuners on single-agent RL, and JumpStart~\cite{jumpstart} ranking agreement against budget in offline RL (it plateaus after 10--25 trials on Atari, and keeps rising on continuous control). Muon has been tested in RL only for LLM post-training, where it collapsed in RLVR~\cite{pion} or helped depending on the estimator~\cite{muonagentic}. POPGym Arcade~\cite{popgym} showed return alone confounds memory evaluation (single-agent). Prior Figgie work is a non-learning simulation~\cite{disilvio} and the CS224R report~\cite{cs224r}. We found no tuner$\times$optimizer factorial for game-solving RL, no ranking-vs-budget curves in multi-agent games, and no test of whether memory-architecture rankings survive tuning against a ground-truth belief. H0, H1 and H2 would each be new for multi-agent games, so the study pays off either way.

\section{Shared interface with the IIG half (agreed with Rich)}
Rich builds and versions a small package, \file{tsinghua-rl/iig/tune/}, which Figgie imports as a pinned git dependency, without a fork: random-search, PROTEIN and CARBS adapters, the shared search-space spec, a Muon/AdamW optimizer factory, budget accounting, Trackinizer SeedSequence seeds, and the ledger plus Trackinizer pre-registration. Both halves give each cell the same total env-step budget $B$, seeds included, and report the same statistics: rank per cell, Kendall $\tau$ between tuned and published rankings, 10 fresh confirmation seeds. $B$ is a placeholder until Rich's PROTEIN pilot sets it. We budget his working assumption, $B=1.5\times10^9$ (150 runs $\times$ 10\,M steps, 3 seeds per screened trial), and E1 scales linearly if it changes. Each game supplies only an env adapter and an eval metric. Rich has accepted the five things Figgie needs: (1) Budget in env steps and \emph{core-hours}, since Figgie is CPU-bound and GPU-hours are the wrong unit. (2) Recurrent policies in the Muon factory: a rule for GRU/minGRU gate matrices (stacked or per gate), with embeddings, heads and biases on AdamW. (3) A noisy, maximised objective (profit is $\pm$4 \$/game per 1000 games): adapters take a noise estimate and a sign, and Figgie reports seat-rotated duplicate games. (4) Asynchronous batches and resume: 12--16 trials run at once per 32-core box, finish out of order and span 12\,h leases, so tuner state must checkpoint. (5) Game-specific hyperparameters (BC lift weight, self-play share, shaping) on top of the shared core (lr, entropy, clip, GAE $\lambda$, KL coefficient).

\section{Experiments}
E0 (before any tuning, $\approx$0 cost). Common-trajectory probes: every network, trained and random-init, run over the same 3000 recorded games, removing the policy confound. A known-answer recall check per core (a cue at tick 0, probed at tick 239) separates ``memory does not help'' from ``cannot learn memory at this length''. MMD added to \file{figgie/rl/train.py} as PPO plus entropy plus a KL term to the previous policy, which is what upstream MMD runs~\cite{rudolph} (\file{mmd.py:469}).

E1, the factorial. 3 tuners (random search, Rudolph's protocol, as control, CARBS, PROTEIN) $\times$ 2 optimizers $\times$ 2 algorithms (PPO, MMD) $\times$ 3 cores (none, minGRU, GRU) $=$ 36 cells, budget $B$ each, multi-fidelity from 250 to 2000 updates (4--33\,M steps). Random search with AdamW is the control. Tuning objective: duplicate profit against the held-out table FV + Bayes + MM. A second tuner replicate on 4 cells measures run-to-run tuner variance.

E2, confirmation. Each cell's winner and the current recipe, 10 fresh seeds $\times$ 2000 updates, on common-trajectory probes and the 6-matchup eval. Approximate exploitability~\cite{timbers}: one PPO best response (20\,M steps, against 3 copies) for each of the 36 winners and 3 current-recipe agents.

Out of scope on Figgie, with reasons. ESCHER and DREAM, as Phase 1 measured 40--185 CPU-h per ESCHER run on the small games alone. A CFR+ floor: no tabular solution exists. We report the best heuristic (AdaptiveMM, $+22.2\pm2.4$ \$/game in our tournament) and the exact posterior for probes, but this is a limitation: neither is an equilibrium bound, and a learned best response only bounds exploitability from below. The transformer core (4$\times$ the cost, degenerate so far) and mini-generals or NLHE move to Jan--Feb.

\section{Pre-registered predictions (filed in the ledger and Trackinizer before E1 launches)}
Our past misses underpredicted the new method (minGRU predicted 43--52\%, observed 58.0\%), so the intervals are wide and skewed toward the treatment. Baselines: current recipe, 3 seeds, none core $+14.6$ \$/game against FV + Bayes + MM.

{\small\setlength{\tabcolsep}{4pt}
\begin{tabular}{@{}p{0.52\textwidth}cc@{}}
\toprule
Quantity (E2, 10 seeds unless stated) & Median & 90\% interval \\
\midrule
P1 Tuning gain, random$\times$AdamW, PPO$\times$none, minus current recipe (\$/game) & $+4$ & $[-2, +15]$ \\
P1b Same, best tuner and optimizer & $+6$ & $[-2, +20]$ \\
P2 Best tuned memory core minus tuned none, common-trajectory probe (pp) & $+2$ & $[-4, +12]$ \\
P3 Same, profit (\$/game) & $0$ & $[-6, +10]$ \\
P4 MMD minus PPO, tuned, same core (\$/game) & $+1$ & $[-6, +8]$ \\
P5 Muon minus AdamW, matched trials, mean over cells (\$/game) & $0$ & $[-6, +8]$ \\
P6 Share of $B$ for PROTEIN (CARBS) to match random search's best & 0.3 (0.4) & $[0.1, >1]$ \\
P7 Kendall $\tau$, 6 arms ranked at 25\% vs 100\% of $B$ & 0.5 & $[-0.1, 1.0]$ \\
P8 Kendall $\tau$, published-recipe vs tuned ranking & 0.3 & $[-0.5, 1.0]$ \\
P9 Best-response gain vs best tuned agent $\div$ vs current PPO-v2 & 0.8 & $[0.4, 1.3]$ \\
\bottomrule
\end{tabular}}

Decision rules. ``Memory helps'' keeps the rule in \file{docs/PREREGISTRATION.md} (mean $\geq$3 points over none), now on common trajectories, 10 seeds, one-sided Mann--Whitney, Holm over the cores. CARBS and PROTEIN each against random search, and Muon against AdamW: sign tests over matched cell pairs (12, 12 and 18), Holm, plus 90\% bootstrap CIs. Transfer (stated, not tested): tuner and optimizer effects have the same sign as in Rich's half, probability 0.6 each. Everything else is exploratory.

\section{Compute}
Training is CPU-bound (Python simulation and opponent forward passes, a 128-wide learner), so we count cores: a 2000-update run costs 5 core-h for none and 7 for minGRU/GRU, measured (\file{docs/COMPUTE\_PLAN\_FINAL.md}).

{\small\setlength{\tabcolsep}{4pt}
\begin{tabular}{@{}lrl@{}}
\toprule
Item & Core-h & Basis \\
\midrule
E1, 36 cells $\times$ $B$ plus 4 replicate cells, incl.\ eval & 11{,}740 & 235 (none) or 325 core-h per cell \\
E2 confirmation, 39 arms $\times$ 10 seeds, probes and eval & 3{,}240 & 5--7 + 2 core-h per run \\
E2 best responses (39 $\times$ 20\,M steps), E0 recall check & 380 & 6 core-h each, 8 runs $\leq$4000 updates \\
Contingency (25\%) & 3{,}840 & reruns, longer runs if recall needs them \\
\textbf{Total} & \textbf{19{,}200} & \\
\bottomrule
\end{tabular}}

Where. A rented 32-core box with a 3090-class GPU is quoted at \$0.20/h on Vast (16 Nov), about \$0.007 per core-hour, so one finished 2000-update run costs about \$0.05 quoted. A lone 8-core 5090 costs about 20$\times$ more per core, with its GPU idle. Total: 19{,}200 core-h $\approx$ 640 box-hours, \textbf{about \$160} (\$128 quoted $\times$ 1.25, the top of the 0.92--1.25 billed-to-quote range measured on 16 Nov, $\approx$1.1 typical). Five boxes finish E1 in about 4 days. The pilot (below) replaces these estimates with measured core-hours per trial. If the cap is lower, we cut in this order: the tuner replicates (to \$148), the GRU core, keeping none vs minGRU (to \$95), then MMD, keeping PPO (to \$50: 12 cells, all 3 tuners $\times$ 2 optimizers at the matched $B$). The matched budget and the 10 seeds go last, since they make the halves comparable. A throughput lever we will try but do not budget on: a C engine with batched opponent inference, in the PufferLib style~\cite{puffer}, gated on bit-exact parity with \file{figgie/engine.py} over 10{,}000 games. (Rich's C port of the small games gained $\approx$150$\times$.)

\section{Milestones}
17--21 Nov: pilot, inside the \$5.5 left of tenant geoff's share: one tuner (PROTEIN, or a stub until \file{iig/tune} exists) on the none core, short runs, one rented 32-core box, $\leq$\$5, pre-registered, results exploratory. It measures core-hours per trial, worker scaling (2/4/8) and \$ per cell. In parallel, E0 on compute-01 and MMD in the trainer. $\approx$23 Nov, midterm presentation: the 3-seed memory table, common-trajectory probes, the pilot and this plan. 23--29 Nov: E1, if the cap is set by 22 Nov. 30 Nov--8 Dec: E2. Each day the cap slips moves both by a day. 9--20 Dec: scoring, figures, poster (week 15, 21--27 Dec). Final report 8 Jan. Jan--Feb (joint paper, outside the course): merge with Rich's half, transformer core, mini-generals or NLHE.

{\scriptsize\begin{multicols}{2}
\begin{thebibliography}{99}\setlength{\itemsep}{-2pt}\renewcommand{\section}[2]{}
\bibitem{rudolph} M.~Rudolph et al. Reevaluating policy gradient methods for IIGs. ICLR 2026, arXiv:2502.08938.
\bibitem{eimer} T.~Eimer, M.~Lindauer, R.~Raileanu. Hyperparameters in RL and how to tune them. ICML 2023, arXiv:2306.01324.
\bibitem{jumpstart} N.~Omi et al. JumpStart your policy learning with lessons from 160,000 training runs. arXiv:2609.13730, 2026.
\bibitem{adkins} J.~Adkins, M.~Bowling et al. Evaluating hyperparameter sensitivity in RL. arXiv:2412.07165, 2024.
\bibitem{pion} Pion. arXiv:2605.19282, 2026.
\bibitem{muonagentic} When does Muon help agentic reinforcement learning? arXiv:2607.16169, 2026.
\bibitem{popgym} Investigating memory in model-free RL with POPGym Arcade. arXiv:2503.01450, 2025.
\bibitem{disilvio} Traders in a strange land: agent-based simulation of Figgie. arXiv:2110.00879, 2021.
\bibitem{cs224r} D.~Li Yang. Reinforcement learning for Figgie. Stanford CS224R project report, 2026.
\bibitem{carbs} A.~Fetterman et al. Tune as you scale (CARBS). arXiv:2306.08055, 2023. PROTEIN: puffer.ai/blog.html.
\bibitem{timbers} F.~Timbers et al. Approximate exploitability. arXiv:2004.09677.
\bibitem{puffer} J.~Suarez. PufferLib. arXiv:2406.12905, 2024.
\end{thebibliography}\end{multicols}}
\end{document}
